\section{Related Work}
\label{sec:related}

\para{Information economics of certification.} Our analytical layer is rooted in the classic signaling tradition. \citet{akerlof1970lemons} shows that markets with asymmetric information collapse without credible signals; \citet{spence1973job} introduces costly signaling; \citet{long2023lemons} derives a closed-form regulator-as-certifier model that maps directly onto a journal serving as a quality stamp. These models are static in volume: the number of agents being screened does not appear as an independent variable. We embed the same Bayesian update inside a capacity-constrained gatekeeper and study its behavior as $\Nsub$ grows, recovering the prior closed-form results in the no-capacity limit.

\para{Empirics of gate-keeping.} \citet{heckman2019tyranny} document that publication in the \emph{Top-Five} economics journals is a powerful but imperfect predictor of tenure, and that the screen captures only about half of the truly influential articles in the field. \citet{kwiek2026topperformers} show that the top decile of researchers by output dominates publishing in top-decile journals across $30$ years of Polish science, and that the gap is qualitative (where) rather than just quantitative (how much). Both studies provide snapshots; neither models how the gap should change with $\Nsub$. We reproduce their qualitative regularities (top-tier captures roughly half of top papers; persistent concentration) as natural cells of the same model.

\para{Volume growth and information per paper.} \citet{bornmann2015growth} estimate that scientific output has been growing at $8$--$9\%$ per year since World War II, doubling every $8$--$13$ years. \citet{fu2021disproportion} use a knowledge-quantification index on $185$M articles to argue that the amount of knowledge contained in the literature grows much more slowly than the literature itself. These bibliometric studies establish the explanatory variable $\Nsub(t)$ but do not connect it to the signaling channel. We provide the missing link: the same volume growth implies signal degradation through the capacity-constraint and load-induced-noise channels.

\para{Peer review noise and dynamics.} The Cortes--Lawrence \neurips~2014 reproducibility experiment found that roughly half of accepted papers would be rejected by a second independent committee, anchoring a plausible reviewer-noise scale. \citet{wang2025peerreview} estimate that $\sim 100$M reviewer-hours per year are now spent globally and document a secondary diffusion-of-ideas function of review. \citet{latona2024lottery} report that LLM-assisted reviewing systematically inflates scores. \citet{mohanty2025crisis} argue that AI-conference reviewing has crossed a sustainability threshold. Our load-induced-noise mechanism gives a quantitative account of why the empirical \iclr signal degrades even though acceptance rate has stayed roughly constant.

\para{Simulation models of peer review.} \textsc{AgentReview}~\citep{jin2024agentreview} simulates a multi-stage \openreview pipeline with LLM agents; \citet{grimaldo2018reputation} compare reputation- and review-based gating in an agent-based model. Both fix the load $\Nsub$ rather than treating it as the key independent variable, and neither couples a reader-side consumption layer or a researcher-level identifiability analysis. Our simulator is intentionally minimal---no LLMs, $<60$s wall-clock to reproduce all experiments---so the dependence on $\Nsub$ is the first-class object of study.

\para{Datasets.} We rely on the open \papercopilot dataset~\citep{yang2025papercopilot} which catalogues submissions, reviews, and scores across major AI conferences (\iclr 2013--2026, \neurips 1987--present, \icml 2013--2025). To our knowledge, this is the first study to use the full \iclr time series to estimate a log-linear decay of standardized posterior contrast in $\Nsub$.

\para{Position.} Building on Akerlof--Spence--Long, we treat gate-keeping as Bayesian signaling. Building on Heckman--Moktan and Bornmann--Mutz, we anchor the empirical regime of interest. Unlike all of these, we make $\Nsub$ the explanatory variable, expand the model to three coupled layers (paper, reader, researcher), and calibrate it directly on a long time series of conference data.
